\section*{Chapter \ref{ch:g3:separating-mixtures} --- Separating Mixtures}

\begin{solution}{exo:g3:separating-mixtures:1}
\omterm{def:g3:separating-mixtures:sieve}{Sieving}: the rice grains are smaller than the peas and fall through a
sieve with the right holes. (Sorting by hand also works, slowly.)
\end{solution}

\begin{solution}{exo:g3:separating-mixtures:2}
The soil sinks slowly and forms a layer at the bottom; the water above
becomes clearer. This is \omterm{def:g3:separating-mixtures:decant}{settling}.
\end{solution}

\begin{solution}{exo:g3:separating-mixtures:3}
The \omterm{def:g3:separating-mixtures:filter}{filtrate} is the liquid that passes through the filter. When coffee
is made, the \omterm{def:g3:separating-mixtures:filter}{filtrate} is the coffee in the jug; the grounds stay in the
paper.
\end{solution}

\begin{solution}{exo:g3:separating-mixtures:4}
Yes. The salt is dissolved, and a filter does not hold back what is
dissolved: it passes through with the water.
\end{solution}

\begin{solution}{exo:g3:separating-mixtures:5}
By \omterm{def:g3:separating-mixtures:evaporate}{evaporating} the water: in the sun, or gently heated, the water goes
into the air and the sugar stays in the dish.
\end{solution}

\begin{solution}{exo:g3:separating-mixtures:6}
The gravel stays in the sieve. The grains of sand are smaller than the
holes, so they fall through.
\end{solution}

\begin{solution}{exo:g3:separating-mixtures:7}
(a) \omterm{def:g3:separating-mixtures:sieve}{Sieving}. (b) \omterm{def:g3:separating-mixtures:filter}{Filtering} (or \omterm{def:g3:separating-mixtures:decant}{settling} and \omterm{def:g3:separating-mixtures:decant}{decanting}). (c) \omterm{def:g3:separating-mixtures:evaporate}{Evaporating}.
\end{solution}

\begin{solution}{exo:g3:separating-mixtures:8}
Four steps: screens, \omterm{def:g3:separating-mixtures:decant}{settling} tanks, sand filters, germs killed. The
\omterm{def:g3:separating-mixtures:decant}{settling} tanks remove the mud.
\end{solution}

\begin{solution}{exo:g3:separating-mixtures:9}
Filter; collect the sand from the paper; evaporate the \omterm{def:g3:separating-mixtures:filter}{filtrate} to get
the salt.
\end{solution}

\begin{solution}{exo:g3:separating-mixtures:10}
No. The salt of sea water is dissolved, and a filter does not hold back
what is dissolved: the \omterm{def:g3:separating-mixtures:filter}{filtrate} is still salty water.
\end{solution}

\begin{solution}{exo:g3:separating-mixtures:11}
Pouring the tea off gently once the leaves have settled (\omterm{def:g3:separating-mixtures:decant}{decanting}), or
pouring it through a strainer or a filter (\omterm{def:g3:separating-mixtures:sieve}{sieving} or \omterm{def:g3:separating-mixtures:filter}{filtering}).
\end{solution}
